\section{Conclusion et perspectives}
\label{sec:conclu}

\subsection{Bilan des réalisations}

Ce projet a permis de concevoir, implémenter et valider un \textbf{accélérateur matériel SHA-256} complet, depuis la spécification algorithmique jusqu'à la vérification fonctionnelle par simulation. Le co-processeur résultant satisfait l'ensemble des objectifs fixés initialement :

\begin{itemize}
    \item \textbf{Conformité FIPS 180-4} : l'implémentation produit des condensés identiques bit à bit aux valeurs de référence publiées par le NIST, validant la fidélité de la transcription algorithmique en logique matérielle. Les 64 constantes $K_t$, le vecteur d'initialisation $H_{init}$, les six fonctions logiques, et la procédure d'accumulation multi-blocs sont tous conformes à la norme.

    \item \textbf{Support multi-blocs} : le co-processeur gère correctement les messages de longueur arbitraire nécessitant le traitement séquentiel de plusieurs blocs de 512 bits. Le mécanisme d'accumulation progressive ($H^{(i)} = H^{(i-1)} + \text{round}(H^{(i-1)}, M_i)$) est validé par le vecteur de test à deux blocs. L'architecture supporte jusqu'à 255 blocs (limité par le champ 8 bits du compteur), soit des messages jusqu'à 16 320 octets dans cette version --- suffisant pour la très grande majorité des applications pratiques.

    \item \textbf{Architecture modulaire} : la décomposition en cinq modules VHDL clairement séparés (Package, Round, Scheduler, FSM, Core) plus le module Memory facilite la compréhension, la maintenance, et la réutilisation. Chaque module a une interface bien définie et peut être testé individuellement. Cette modularité permet également le remplacement d'un module sans impacter les autres (par exemple, remplacer le module Round par une version optimisée carry-save).

    \item \textbf{Validation fonctionnelle} : trois vecteurs de test couvrant les cas mono-bloc et multi-blocs, avec un banc de test auto-vérifiant qui produit un verdict PASS/FAIL sans intervention manuelle. La correspondance parfaite des 256 bits de chaque condensé avec les références démontre une implémentation sans erreur.

    \item \textbf{Performance démontrée} : un débit théorique de 602 Mbit/s à 100 MHz (764 Mbit/s asymptotique), représentant un gain de 24 à 48 fois par rapport à un processeur embarqué ARM Cortex-M4. Ce gain justifie pleinement le qualificatif d'accélérateur matériel.
\end{itemize}

\subsection{Compétences développées}

Au-delà du résultat technique, ce projet a permis de développer et d'approfondir un ensemble de compétences en conception numérique :

\paragraph{Conception RTL} La traduction d'un algorithme séquentiel (décrit par des boucles et des variables en pseudo-code) en une architecture matérielle parallèle (décrite par des registres, des chemins combinatoires, et une machine à états) est un exercice intellectuel fondamentalement différent de la programmation logicielle. Il nécessite de penser en termes de cycles d'horloge, de pipelines, et de parallélisme physique plutôt qu'en termes d'instructions séquentielles.

\paragraph{Gestion du pipeline} La maîtrise de la latence mémoire et du séquencement des signaux de contrôle est une compétence critique en conception de circuits intégrés. Les bugs de pipeline (section~\ref{sec:bugs}) ont fourni une expérience concrète de diagnostic et de correction de ces problèmes subtils.

\paragraph{Vérification fonctionnelle} La conception d'un testbench auto-vérifiant avec des vecteurs de test pertinents, le diagnostic d'erreurs par analyse de chronogrammes, et la compréhension des limites de la simulation fonctionnelle par rapport à la simulation temporelle constituent des compétences directement transposables à un contexte industriel.

\paragraph{Optimisation architecturale} L'analyse comparative des architectures (déroulée vs pipelinée vs itérative) et la justification du choix par des critères quantitatifs (surface, débit, fréquence) illustrent la démarche d'ingénierie systématique nécessaire à la conception de systèmes performants sous contraintes.

\subsection{Perspectives d'amélioration}

Plusieurs axes d'amélioration permettraient de transformer ce prototype académique en un composant IP (Intellectual Property) de qualité industrielle :

\paragraph{1. Interface AXI-Lite pour intégration SoC Zynq} L'interface mémoire actuelle (signaux discrets) est fonctionnelle mais non standard. L'ajout d'un wrapper AXI-Lite permettrait une intégration directe dans un système Zynq-7000 via le bus AXI du processeur ARM Cortex-A9 embarqué. Ce wrapper convertirait les transactions AXI (adresse, donnée, réponse, handshake) en signaux internes du co-processeur, avec un registre de contrôle pour le déclenchement et un registre de statut pour la notification de fin. Le pilote Linux correspondant serait trivial (memory-mapped I/O).

\paragraph{2. Pré-calcul $K_t + W_t$ pour augmenter la Fmax} Le chemin critique actuel traverse cinq additions de 32 bits en série dans le module Round. En pré-calculant la somme $K_t + W_t$ dans un registre dédié un cycle avant son utilisation, on réduit la chaîne critique à quatre additions. Cette optimisation, classique dans la littérature~\cite{dadda}, offre un gain de Fmax estimé à 20--25\% (passage de 150 MHz à 180--190 MHz sur Artix-7) au prix d'un cycle de latence supplémentaire par bloc et de 32 bascules additionnelles.

\paragraph{3. Architecture multi-coeurs} Pour les applications nécessitant un débit supérieur (minage de crypto-monnaie, serveur TLS haute performance), plusieurs instances du co-processeur peuvent être instanciées en parallèle, chacune traitant un message différent. Avec $N$ coeurs, le débit est multiplié par $N$ de manière linéaire (les coeurs sont indépendants). Un Artix-7 XC7A100T (53\,000 LUT) pourrait accueillir environ 25 coeurs, offrant un débit agrégé de 15 Gbit/s.

\paragraph{4. Support HMAC-SHA256} Le protocole HMAC (Hash-based Message Authentication Code) utilise deux applications consécutives de SHA-256 avec une clé secrète. L'ajout du support HMAC nécessiterait un registre de clé dédié et une logique de contrôle pour enchaîner les deux passes de hachage (inner hash et outer hash) automatiquement. Le coût en surface est minimal (un registre de 256 bits et une extension de la FSM), mais le bénéfice est considérable car HMAC-SHA256 est omniprésent dans TLS, IPsec, et OAuth.

\paragraph{5. Synthèse et validation timing sur FPGA réel} L'étape logique suivante serait la synthèse complète du design sur un FPGA cible (Artix-7 ou Zynq-7000), permettant de mesurer la fréquence maximale réelle post-placement-routage, la consommation de ressources exacte (LUT, FF, BRAM, DSP), et la puissance dissipée. Cette étape révélerait les éventuels problèmes de timing (slack négatif sur le chemin critique) et nécessiterait potentiellement des contraintes de placement ou des modifications RTL pour les résoudre.

\paragraph{6. Extension du compteur de blocs} Le champ \texttt{num\_blocks} actuel est limité à 8 bits (255 blocs maximum = 16 320 octets). Pour supporter des messages plus longs (fichiers, flux réseau), ce champ devrait être étendu à 16 ou 32 bits. Cette modification est triviale en termes de RTL (élargir le signal et le comparateur) mais nécessite une mise à jour de la carte mémoire et du protocole hôte.

\subsection{Réflexion finale}

Ce projet illustre concrètement la puissance de la conception matérielle dédiée pour l'accélération de fonctions critiques. Là où un processeur généraliste exécute une séquence d'instructions une à une, un circuit dédié exploite le parallélisme inhérent à l'algorithme pour atteindre des performances inaccessibles au logiciel. Le gain de 24--48$\times$ obtenu avec un circuit occupant moins de 4\% d'un FPGA modeste démontre que l'accélération matérielle n'est pas réservée aux grandes entreprises ou aux projets à gros budget : avec les outils de développement FPGA modernes (Vivado, GHDL, cocotb), un ingénieur individuel peut concevoir un co-processeur cryptographique compétitif en quelques semaines.

La tendance actuelle de l'industrie vers les architectures hétérogènes (CPU + GPU + FPGA + accélérateurs dédiés) confirme la pertinence de cette approche. Les SoC modernes (Xilinx Zynq, Intel Agilex, Apple M-series) intègrent tous des accélérateurs matériels pour les fonctions critiques (crypto, vidéo, IA), et la maîtrise de leur conception constitue une compétence de plus en plus valorisée dans l'industrie du semi-conducteur.

\newpage
\begin{thebibliography}{9}
\bibitem{fips180} NIST, \emph{Secure Hash Standard (SHS)}, FIPS PUB 180-4, Information Technology Laboratory, National Institute of Standards and Technology, Gaithersburg, MD, August 2015.
\bibitem{mcloone} M. McLoone, J. V. McCanny, ``High Performance Single-Chip FPGA Rijndael Algorithm Implementations and SHA-384/512 Architectures'', in \emph{Proc. IEEE International Conference on Field-Programmable Logic (FPL)}, pp. 256--265, 2002.
\bibitem{dadda} L. Dadda, M. Macchetti, J. Owen, ``The Design of a High Speed ASIC Unit for the Hash Function SHA-256 (384, 512)'', in \emph{Proc. ACM Great Lakes Symposium on VLSI (GLSVLSI)}, pp. 25--28, 2004.
\bibitem{icact2019} M. Dowenbergh, A. Al-Anbuky, S. Garg, ``SHA-256 Hardware Acceleration on FPGA-CPU Heterogeneous Cluster'', in \emph{Proc. 21st International Conference on Advanced Communication Technology (ICACT)}, pp. 412--418, 2019.
\end{thebibliography}

